\subsection{A local gradient estimate}

The boundedness established above does not yet permit ordinary Schauder
estimates for the second equation: its source contains the square of the
Hessian of $f$.  We first prove a scalar estimate that also controls the
local mass of that square.  Throughout this subsection $D_x$ denotes
coordinate differentiation, whereas $\nabla$ and $\Hess$ refer to the
background Riemannian metric.  The balls used in coordinate estimates are
Euclidean coordinate balls; uniformly equivalent metric balls give the
same conclusions after changing constants.

\begin{lemma}[Local gradient and Hessian estimates]\label{g:gradient}
Consider three-dimensional coordinate patches with a fixed positive lower
bound for their size, uniform ellipticity and smooth bounds for their
metrics, and, when present, uniform smooth geometry of one boundary face.
Let $z$ be a smooth solution, including up to that face, of
\begin{equation}\label{g:axial-equation}
 A_\chi:\Hess z=G_0,\qquad
 A_\chi=I-(1-\chi)e\otimes e,\qquad
 e=\frac{\nabla z}{|\nabla z|},\qquad
 0<\chi_0\leq\chi\leq1.
\end{equation}
At a zero of $\nabla z$ any measurable unit axis for which the equation
holds is allowed.  Assume
$|\nabla z|\leq L_0$ and $|G_0|\leq Q_0$; on a boundary face assume that
$z$ is constant.  Only the displayed measurable bounds on $\chi$ and
$G_0$ are used.  There are $\alpha\in(0,1)$, $C<\infty$, and a uniform
smaller patch size, depending on $\chi_0$ and the geometric bounds, such
that
\begin{equation}\label{g:gradient-estimates}
 \|D_xz\|_{C^{0,\alpha}}\leq C(L_0+Q_0),\qquad
 \int_{B_r}|\Hess z|^2\,\dd V_g
       \leq C(L_0+Q_0)^2r^{1+2\alpha}.
\end{equation}
The first norm is on the smaller patch.  The second estimate holds for
all sufficiently small balls centered in it, with intersection with the
domain understood at a boundary.  The same estimate holds for $D_x^2z$
and coordinate volume.  These are a priori estimates for the stated
smooth solutions, with constants independent of their higher derivatives.
\end{lemma}

\begin{proof}
We organize the proof so that neither a derivative nor a continuity
modulus of $\chi$ enters.

\paragraph{Reflection and a Hessian estimate near exponent two.}
At a constant Dirichlet face subtract the boundary value and use boundary
normal coordinates.  Reflect $z$ oddly and the metric isometrically.
Tangential first derivatives vanish on the face and the normal first
derivatives of the odd extension agree.  Thus the extended $z$ belongs
locally to $W^{2,p}$ for every finite $p$ for each individual smooth
solution; its second derivative has no surface atom.  The reflected metric
is Lipschitz and smooth on each closed side, with uniform bounds there.
Equation~\eqref{g:axial-equation} holds almost everywhere on both sides,
with the reflected axis and the odd reflected source.  A linear coordinate
change at the center makes the metric Euclidean there.  On sufficiently
small rescaled balls the coordinate principal matrix, denoted by $A$, obeys
\begin{equation}\label{g:cordes-defect}
 \|I-A\|_{\mathrm F}\leq1-\chi_0/2<1.
\end{equation}
Indeed its defect at a point in an orthonormal frame is precisely
$1-\chi$, and the coordinate metric error can be made smaller than
$\chi_0/2$.  This statement also holds in a reflected patch.

The strict-defect Hessian absorption below is a Cordes-type estimate;
for the linear framework see \cite[Section~2, Lemma~1, and Section~2.1,
Theorems~2--3]{SmearsSuli2013}.  The reflected near-two estimate and the
subsequent nonlinear gradient and Morrey bounds are developed here.

For a compactly supported function on $\R^3$, the operator
$D_x^2\Delta^{-1}$ from a scalar source to the full Hessian array has
$L^2$ norm one: its Fourier multiplier has squared Frobenius norm
$\sum_{i,j}\xi_i^2\xi_j^2/|\xi|^4=1$.  Its $L^p$ norm is bounded for
$1<p<\infty$, and interpolation makes that norm arbitrarily close to one
when $p$ is sufficiently close to two.  Absorbing
$(I-A):D_x^2v$ using~\eqref{g:cordes-defect} therefore gives, for some
$2<p_0<3$,
\begin{equation}\label{g:cordes}
 \|v\|_{W^{2,p_0}(B_{1/2})}
 \leq C\bigl(\|v\|_{L^{p_0}(B_1)}
        +\|A_\chi:\Hess v\|_{L^{p_0}(B_1)}\bigr).
\end{equation}
Here and below a fixed smaller pair of concentric balls may replace the
displayed pair.  To localize, apply the compact-support estimate to a
cutoff times $v$ and absorb first derivative terms by
$\|D_xv\|_p\leq\eta\|D_x^2v\|_p+C_\eta\|v\|_p$.
The rescaled Christoffel terms are handled by the same inequality.
Using nested cutoffs gives an arbitrarily small multiple of the larger
Hessian norm, with a polynomial cost in the inverse gap; iteration on
nested balls absorbs that term.  The proof also gives the exponent-two
version of~\eqref{g:cordes}.  It applies to the piecewise smooth reflected
functions just described.

\paragraph{A flux identity and a drop or affine alternative.}
Normalize the gradient bound to one on a rescaled unit ball and make the
source and geometric errors at most $\delta$.  Geometric errors here
include the metric deviation from its constant value, its piecewise first
derivatives, the smooth-side curvature, and the scaled second fundamental
form of a reflection face.  Put
$P=\nabla z$, $H_z=\Hess z$, and $s=1-|P|^2\geq0$.  The equation gives
$\Delta z=G_0+(1-\chi)H_z(e,e)$ and therefore the exact identity
\begin{equation}\label{g:flux-cancellation}
 A_\chi\nabla\frac{|P|^2}{2}-G_0P
          =(H_z-(\Delta z)g)P.
\end{equation}
This identity remains valid when $P=0$.  Since
$1-\chi\leq1-\chi_0<1$, Young's inequality, with a fixed parameter
depending on $\chi_0$, gives
\begin{equation}\label{g:hessian-coercivity}
 |H_z|^2-(\Delta z)^2\geq c|H_z|^2-CG_0^2.
\end{equation}
For example, choose $\eta>0$ with
$(1+\eta)(1-\chi_0)^2<1$ and use
$(G_0+(1-\chi)H_z(e,e))^2
 \leq(1+\eta)(1-\chi_0)^2|H_z|^2+C_\eta G_0^2$.
The divergence of the right side of~\eqref{g:flux-cancellation} is
\[
 |H_z|^2-(\Delta z)^2+\Ric(P,P).
\]
Consequently, on the smooth sides,
\begin{equation}\label{g:gradient-deficit}
 -\divg\bigl(A_\chi\nabla s+2G_0P\bigr)
             \geq2c|H_z|^2-C\delta.
\end{equation}
The differentiation is of the right side of the exact identity;
it does not require differentiating $\chi$ or $G_0$ separately.

There is also a controlled distributional version across a reflected
face.  There $P=(\partial_\nu z)\nu$, and the normal component of the
right side of~\eqref{g:flux-cancellation} is
\[
 -\bigl(\tr_T H_z\bigr)\partial_\nu z.
\]
The tangential Hessian of a function constant on the face is, up to the
normal-sign convention, its normal derivative times the second fundamental
form.  Thus the flux jump has magnitude at most $C\delta$, irrespective
of the second normal derivative of $z$.  Incorporate the volume density
in coordinate divergence.  A bounded plane density $q(x')$ is the
divergence of the bounded field
$q(x')\mathbf1_{\{x_3>0\}}\partial_{x_3}$ locally.  Hence the plane
error in~\eqref{g:gradient-deficit} is another divergence-field error of
size $C\delta$.  All subsequent weak inequalities have bounded uniformly
elliptic principal coefficients and scalar and vector errors tending to
zero with $\delta$.

Consider a sequence with $\delta\to0$ and
$\inf_{B_{1/2}}s\to0$.  Dropping the favorable Hessian term,
the inhomogeneous weak Harnack inequality for nonnegative divergence
supersolutions gives $s\to0$ in measure on interior smaller balls
\cite{GilbargTrudinger2001}.  Fixed radii in this argument can be chosen
so as to include any prescribed ball compactly contained in $B_{1/2}$.
We also have $D_xs\to0$ in $L^1$ locally there.  To see the latter claim,
test the supersolution inequality by $\eta^2(b-s)_+$, where $b>0$ and
$\eta$ is a compactly supported cutoff.  Ellipticity and Young's
inequality give
\begin{equation}\label{g:low-level}
 \int_{\{s<b\}}\eta^2|D_xs|^2
                  \leq C_\eta b^2+o(1).
\end{equation}
The gradient bound and~\eqref{g:cordes} already give a uniform local
$L^2$ bound for $D_xs$.  On $\{s\geq b\}$, Cauchy's inequality and
convergence in measure therefore give convergence of its $L^1$ integral
to zero.  On $\{s<b\}$ use~\eqref{g:low-level}, then let $b\downarrow0$.
Testing the full inequality~\eqref{g:gradient-deficit} by a fixed
nonnegative cutoff now yields $H_z\to0$ locally in $L^2$.  Since the
scaled Christoffel symbols tend to zero, $D_x^2z\to0$ there as well.
After subtracting constants, Poincar\'e's inequality makes the gradients
converge in $L^2$ to a constant vector.  Its length is one because
$s\to0$ in measure.  The gradient bounds give local uniform convergence
of a subsequence to the corresponding affine function.

It follows that for every sufficiently small prescribed $b_*>0$ there
are fixed numbers $0<r_*<k_*<1$ and an error threshold such that one of
the following holds:
\begin{equation}\label{g:drop-or-flat}
 \begin{split}
 &\sup_{B_{r_*}}|\nabla z|\leq k_*;\qquad\text{or}\\
 &\sup_{B_{r_*}}|z-L|\leq b_*r_*,\qquad
                  \tfrac12\leq|D_xL|\leq2
                  \quad\text{for an affine }L.
 \end{split}
\end{equation}
Indeed, fix $r_*<1/4$.  Failure for every $k_*<1$ and every sufficiently
small error threshold would give a sequence with gradients approaching
length one somewhere in $B_{r_*}$ and no such affine approximation.
The preceding compactness conclusion contradicts this.  Increasing
$k_*$ if necessary ensures $r_*<k_*$.

\paragraph{Improving a nonzero affine approximation.}
Rescale the affine branch to a unit ball.  Suppose
\[
 \sup_{B_1}|z-L_1|\leq b,\qquad
 \tfrac14\leq|D_xL_1|\leq4,
\]
the gradient has a fixed bound, and the forcing and scaled geometric
errors are at most $b\delta$.  Set $Y=(z-L_1)/b$.  The Christoffel
contribution of the affine $L_1$ is included in these errors, so
$|A_\chi:\Hess Y|\leq C\delta$.  Equation~\eqref{g:cordes} gives a
uniform interior $W^{2,p_0}$ bound for $Y$.  Let
$e_0=D_xL_1/|D_xL_1|$ be the fixed Euclidean axis.  On an interior ball
the difference between the actual principal matrix and
$I-(1-\chi)e_0\otimes e_0$ tends to zero in each finite $L^q$ norm as
$b\to0$, uniformly in the displayed slope range.  Indeed
$D_xz=D_xL_1+bD_xY$, the last term tends to zero in measure, and the
coefficient difference is bounded.  The small metric error has the same
property.  The set where $D_xz=0$ causes no problem: it is contained in
the set where $b|D_xY|\geq1/4$.
Choose $q$ by $1/q=1/2-1/p_0$.  H\"older's inequality then gives
\begin{equation}\label{g:axis-replacement}
 \| (\Delta_{\perp_0}+\chi\partial_{e_0e_0})Y\|_{L^2(B_{3/4})}
                 \leq C\delta+o_{b\to0}(1).
\end{equation}
The coefficient $\chi$ remains measurable and spatially variable.
The smallness here is uniform for all $b$ below one fixed threshold.
Indeed Sobolev embedding bounds $D_xY$ in
$L^{p^*}$, where $p^*=3p_0/(3-p_0)$.  Splitting according to
$b|D_xY|\leq1$ gives
\[
 \|B_\chi-A_0\|_{L^q}
       \leq Cb^\sigma+Cb\delta,\qquad
 \sigma=\min\{1,p^*/q\}>0,
\]
where $A_0=I-(1-\chi)e_0\otimes e_0$ and $B_\chi$ is the actual
coordinate principal matrix.  This error need only be below the fixed
tolerance for the normalized correction $Y$, not below $b\delta$.
To track all coordinate terms, write $A_0=I-(1-\chi)e_0\otimes e_0$
and let $B_\chi$ be the actual coordinate principal matrix.  With
$q_L=D_xL_1$, the equation used here is
\[
 A_0:D_x^2Y=(A_0-B_\chi):D_x^2Y
       +\frac{G_0}{b}+\frac{B_\chi:\Gamma q_L}{b}
       +B_\chi:\Gamma D_xY.
\]
The notation $B_\chi:\Gamma q$ means
$B_\chi^{ij}\Gamma_{ij}^{k}q_k$.  The last term is bounded in $L^2$
by $Cb\delta$, and the other source and connection terms by $C\delta$.

For completeness, the fixed-axis inhomogeneous equation on $B_{3/4}$
with zero Dirichlet data has a unique $W^{2,2}\cap W^{1,2}_0$ solution.
Write it as
$\Delta v=f+(1-\chi)\partial_{e_0e_0}v$ and iterate the Dirichlet
Laplacian inverse in the norm $\|\Delta v\|_2$.  The convex-ball Hessian
inequality
\[
 \|D_x^2v\|_2\leq\|\Delta v\|_2
\]
makes the contraction factor at most $1-\chi_0$.  Subtract this solution
for the source in~\eqref{g:axis-replacement}.  It has small $W^{2,2}$
norm, hence small uniform norm in dimension three.  The remainder $Y_0$
satisfies
\begin{equation}\label{g:fixed-axis}
 \Delta_{\perp_0}Y_0+\chi\partial_{e_0e_0}Y_0=0
\end{equation}
and has a uniform $W^{2,2}$ norm.

This fixed-axis equation has an interior $C^{1,\alpha_1}$ estimate even
for measurable $\chi$.  To prove it, put $v_1=\partial_{e_0}Y_0$.
Distributional differentiation in the distinguished direction gives
\begin{equation}\label{g:directional-divergence}
 \Delta_{\perp_0}v_1+
               \partial_{e_0}(\chi\partial_{e_0}v_1)=0.
\end{equation}
Here $v_1\in W^{1,2}$ and $\chi\partial_{e_0}v_1\in L^2$; the
formula makes no assertion that $\partial_{e_0}\chi$ exists as a
function.  De Giorgi's estimate makes $v_1$ H\"older continuous, with
some $0<\alpha_1<1$.  Caccioppoli applied after subtracting its value at
the center gives, on smaller balls,
\[
 \int_{B_r}|D_xv_1|^2\leq Cr^{1+2\alpha_1}.
\]
Since
$\Delta Y_0=(1-\chi)\partial_{e_0}v_1$, its Poisson source $q_0$
satisfies
\[
 \int_{B_r}|q_0|\leq Cr^{2+\alpha_1}.
\]
Localize this source in a fixed smaller ball and take its Newton
potential.  The gradient kernel and its derivative have sizes
$C|x-y|^{-2}$ and $C|x-y|^{-3}$.  For two points a distance $h$ apart,
the near and far annuli therefore bound the gradient difference by
\[
 C\sum_{r\lesssim h}r^{\alpha_1}
       +Ch\sum_{r\gtrsim h}r^{\alpha_1-1}
       \leq Ch^{\alpha_1},
\]
where the sums are over dyadic radii within the fixed ball.
The difference between $Y_0$ and this potential is harmonic locally and
has uniform interior bounds.  This proves the asserted estimate for
every component of $D_xY_0$.

Choose $0<\alpha_2<\alpha_1$, then a sufficiently small fixed
$\lambda_0\in(0,1)$, and finally $\delta$ and $b$ sufficiently small.
The affine approximation of $Y_0$ at the center, together with the small
uniform correction, produces an affine $L_2$ with
\begin{equation}\label{g:affine-improvement}
 \sup_{B_{\lambda_0}}|z-L_2|
        \leq b\lambda_0^{1+\alpha_2},\qquad
 |D_xL_2-D_xL_1|\leq Cb.
\end{equation}
More explicitly choose $\lambda_0$ so that the $C^{1,\alpha_1}$
remainder is at most
$\tfrac12\lambda_0^{1+\alpha_2}$, then make the correction smaller
than the same quantity.  This also explains the order of the choices.

\paragraph{Iteration and the Hessian mass bound.}
Choose $b_*$ small enough for~\eqref{g:affine-improvement} and so that
\[
 \frac{Cb_*}{1-\lambda_0^{\alpha_2}}<\frac14.
\]
Then fix the constants in~\eqref{g:drop-or-flat}.  If $L_0+Q_0=0$ the
conclusion is immediate.  Otherwise divide the dependent variable by
$L_0+Q_0$ and start on a sufficiently small uniform spatial scale.
At each occurrence of the gradient-drop branch, divide the spatial
scale by $r_*^{-1}$ and the gradient normalization by $k_*^{-1}$;
equivalently replace $z(y)$ by $z(r_*y)/(r_*k_*)$, up to an additive
constant.  The normalized source is multiplied by $r_*/k_*<1$.
Geometric errors also decrease.  If this branch persists, the constant
approximants give the required approximation with every
$\alpha\leq\log k_*/\log r_*$.

If the affine branch occurs, rescale its ball and iterate
\eqref{g:affine-improvement} with errors
$b_j=b_*\lambda_0^{j\alpha_2}$.  At each step divide the dependent
variable by the spatial ratio, preserving the gradient bound, and
retain the full affine slope.  The total slope drift is less than
$1/4$, so all slopes remain in the range required for the comparison.
The forcing and geometric errors shrink as $\lambda_0^j$, whereas
the required tolerance is $b_*\lambda_0^{j\alpha_2}\delta$.
Because $\alpha_2<1$, an initial scale that enforces this relative
smallness at entry enforces it at every later step.  The same initial
choice works after any number of drops, since $r_*/k_*<1$.

Thus at each center and all small radii there is an affine approximant
with error at most $C(L_0+Q_0)r^{1+\alpha}$, where
\[
 0<\alpha\leq
       \min\{\alpha_2,\log k_*/\log r_*\}<1.
\]
The ratios of successive scales are fixed and the slopes are bounded.
Comparing approximants at successive scales shows convergence of their
slopes; comparing approximants on overlapping balls bounds the difference
of limiting slopes by $C(L_0+Q_0)|x-y|^\alpha$.  For a smooth $z$ each
limiting slope is $D_xz$ at the center.  This gives the first estimate
in~\eqref{g:gradient-estimates}.  All previous arguments apply to the
reflected Lipschitz metric, since its only additional flux error was
already included.

Finally subtract the affine approximant on $B_{2r}$ and apply the
exponent-two estimate~\eqref{g:cordes} at scale $r$.  Its value error
contributes $C(L_0+Q_0)^2r^{1+2\alpha}$ to the Hessian-square integral.
The bounded forcing and the Christoffel term acting on the bounded
affine slope contribute at most $C(L_0+Q_0)^2r^3$.  Since $\alpha<1$,
this is bounded by the same claimed power for small $r$.  Adding back
the Christoffel term proves both versions of the Hessian estimate.
\end{proof}

\subsection{Regularity of the coupled equations}

\begin{proposition}[A priori bootstrap at fixed parameters]
\label{g:regularity}
Fix the prepared data, the selected reduced domain and collars, and all
deformation parameters, including $n$ and $\eps$.  Consider smooth
solutions of the homotopy equations on sufficiently large finite
truncations, on the side $t\geq-\eps$ of the floor.  The bounds for
$f,Z,t,|\nabla f|$ established above imply uniform local classical
estimates, including at the inner and outer boundary faces, through each
prescribed finite order.  Constants are uniform over the homotopy and
the truncations and may depend arbitrarily on the fixed deformation
parameters.  In particular they are not asserted to be uniform as
$n\to\infty$.
\end{proposition}

\begin{proof}
Use uniform patches on the compact core and on the end, including
boundary normal patches.  Boundary components carrying different
conditions are separated; the patches are chosen to meet at most one
face.  At every stage the trace equation is
\begin{equation}\label{g:trace-normalized}
 A_\chi:\Hess f
     =\frac{\sqrt D}{l}\bigl(F-\tr_{A_\chi}K\bigr).
\end{equation}
On the bounded variable range, $l/\sqrt D$ is bounded below and
$\chi$ has a positive lower bound at fixed parameters.  The right side
is bounded, and the face values of $f$ are componentwise constant.
Lemma~\ref{g:gradient} consequently gives a uniform H\"older bound for
$D_xf$ and
\begin{equation}\label{g:f-hessian-measure}
 \int_{B_r}|D_x^2f|^2\,\dd x\leq Cr^{1+2\alpha}.
\end{equation}

\paragraph{The second equation as an equation with a controlled measure.}
Write its coordinate divergence form as
\begin{equation}\label{g:Z-divergence}
 \operatorname{div}_x(a_0D_xZ+b_0)=e_0,
 \qquad |e_0|\leq C(1+|D_xZ|^2+|D_x^2f|^2).
\end{equation}
Here $a_0$ includes the volume density and the principal tensor
$4uA_\chi$; it is bounded, symmetric and uniformly elliptic.  The
bounded flux summand $b_0$ includes the prescribed homotopy drift.
The expression for $t$ is smooth in $(x,Z,D_xf)$ on the bounded range.
Its first derivatives are bounded by
$C(1+|D_xZ|+|D_x^2f|)$, and the tensor expression for $\cT$ is
smooth and quadratic in these derivatives.  These facts prove the
displayed growth bound, including where $D_xf=0$.

For a flux face, flatten the boundary and reflect $Z$ evenly.  For the
separate outer Dirichlet face, reflect $Z$ oddly after subtracting its
constant value.  If $S_0$ is Euclidean reflection in the flattened
plane, reflect $a_0$ as $S_0a_0S_0$ in both cases.  Reflect $b_0$ as
$S_0b_0$ in the even case and as $-S_0b_0$ in the odd case.
For the even extension the bounded prescribed conormal flux gives a
bounded plane density in the divergence.  For the odd extension the
normal component of the total flux agrees on the two sides, so there
is no plane source.  The smooth solution on each closed side justifies
these distributional formulas directly.  By~\eqref{g:f-hessian-measure},
the extended equation has a signed measure right side obeying
\begin{equation}\label{g:measure-growth}
 |e_0|\leq C(1+|D_xZ|^2)\,\dd x+\mu_0,
 \qquad \mu_0(B_r)\leq Cr^{1+\beta},\qquad
 0<\beta\leq\min\{2\alpha,1\},
\end{equation}
where $\mu_0$ is a positive measure.  The upper restriction
$\beta\leq1$ includes the $O(r^2)$ mass of a bounded plane density.

\paragraph{An exponential transformation and oscillation decay.}
Let $\mathcal L_0=-\operatorname{div}_x(a_0D_x)$ and
$V_\pm=e^{\pm kZ}$.  Since $Z$ is bounded, these functions and their
reciprocals have fixed bounds for any fixed $k$.  The product rule in
\eqref{g:Z-divergence} gives
\begin{align}
 \mathcal L_0 V_\pm
  ={}&\operatorname{div}_x(\pm kV_\pm b_0)
       \mp kV_\pm e_0\notag\\
    &-k^2V_\pm(a_0D_xZ+b_0)\cdot D_xZ
 \leq\operatorname{div}_x F_\pm+C(\dd x+\mu_0),
 \qquad |F_\pm|\leq C.\label{g:exponential}
\end{align}
Indeed the negative term containing $k^2a_0$ absorbs the
$Ck|D_xZ|^2$ term from~\eqref{g:measure-growth} when $k$ is fixed
sufficiently large; Young's inequality absorbs the bounded-drift cross
term.  Products with the plane measure are legitimate because $Z$ is
continuous across the reflecting plane.

On a ball $B_h$ in the extended patch solve the zero Dirichlet problem
\[
 \mathcal L_0w_\pm
       =\operatorname{div}_x F_\pm+C(\dd x+\mu_0).
\]
Its solution satisfies
\begin{equation}\label{g:green-correction}
 \|w_\pm\|_\infty\leq\varepsilon_h:=C(h+h^\beta).
\end{equation}
For the bounded divergence field, the scaled vector-source estimate
with any exponent greater than three gives the term $Ch$.  For the
positive measure, the Dirichlet Green function of a bounded measurable
uniformly elliptic divergence operator in dimension three has upper
bound $C|x-y|^{-1}$ \cite{LittmanStampacchiaWeinberger1963}.  If necessary
extend coefficients to a doubled ball before using that bound.  The
growth in~\eqref{g:measure-growth} yields
\[
 \sup_x\int_{B_h}|x-y|^{-1}\,\dd\mu_0(y)
 \leq C\sum_{j\geq0}(2^{-j}h)^\beta\leq Ch^\beta.
\]
The volume source has the smaller bound $Ch^2$.  Existence in
$W^{1,2}_0$ may be seen by smoothing the restricted positive measure
with its uniform ball-growth bound, using these uniform sup estimates,
and passing to a limit.  Testing the smooth problems by their solutions
bounds the energies, since the measure term is at most the sup norm
times its total mass and the vector term is controlled by Cauchy's
inequality.  This also justifies the correction for a measure source.

Equation~\eqref{g:exponential} now shows that
\[
 S_\pm=\sup_{B_h}V_\pm-V_\pm+w_\pm+\varepsilon_h
\]
are nonnegative weak supersolutions of $\mathcal L_0$.  Write
$M_h=\sup_{B_h}Z$ and $m_h=\inf_{B_h}Z$.  On the bounded range of
$Z$, the positive exponential deficit is uniformly comparable to
$M_h-Z$, and the negative exponential deficit to $Z-m_h$.
At each point their ordinary counterparts add to $M_h-m_h$.
Therefore one of the two exponential deficits is at least a fixed
multiple of $M_h-m_h$ on at least half of $B_{h/2}$.
Weak Harnack applied to that $S_\pm$, and then
\eqref{g:green-correction}, bounds the corresponding deficit below
throughout $B_{h/4}$ by
$c(M_h-m_h)-C\varepsilon_h$.  In the positive case this lowers the
upper value of $Z$; in the negative case it raises the lower value.
Consequently
\begin{equation}\label{g:oscillation-decay}
 \operatorname{osc}_{B_{h/4}}Z
       \leq(1-c)\operatorname{osc}_{B_h}Z+C(h+h^\beta).
\end{equation}
Iteration gives a uniform $C^{0,\theta}$ bound for some $\theta>0$.
The same proof in the reflected charts gives the estimate to both kinds
of boundary face.

\paragraph{Absorption of the quadratic gradient term and bootstrap.}
Now $Z$ and $D_xf$ are H\"older continuous, and all coefficients of
\eqref{g:trace-normalized} are smooth functions of those variables and
$x$.  Interior and Dirichlet Schauder estimates give a uniform
$C^{2,\alpha'}$ bound for $f$, for some $\alpha'>0$.
Expand~\eqref{g:Z-divergence} on the original, unreflected patches.
The principal coefficients are H\"older continuous, and the right side
is bounded in absolute value by $C(1+|D_xZ|^2)$, since $D_x^2f$ is
now bounded.  On an inner face $df$ is normal; hence $A_\chi$ preserves
the normal direction, and its positive normal eigenvalue converts the
prescribed conormal flux into
\[
 \partial_\nu Z=\mathcal G(x,Z,f,D_xf),
\]
where $\mathcal G$ is smooth on the bounded range.  The outer condition
is Dirichlet.  There is no junction of these conditions within a patch.

Fix an exponent $p>3$.  The local linear $W^{2,p}$ estimate for these
H\"older principal coefficients follows by freezing coefficients and
absorption, using the usual Dirichlet or normal-derivative estimate at
a smooth face \cite{AgmonDouglisNirenberg1959,GilbargTrudinger2001}.
On nested uniform patches $U\Subset U'$ it gives
\begin{equation}\label{g:W2p-preabsorption}
 \|Z\|_{W^{2,p}(U)}
       \leq C\bigl(1+\|D_xZ\|_{L^{2p}(U')}^2
                         +\|Z\|_{W^{1,p}(U')}\bigr).
\end{equation}
For the boundary term, extend the smooth composition $\mathcal G$
into a collar.  Its $W^{1,p}$ norm is bounded by
$C(1+\|Z\|_{W^{1,p}})$ because $f$ is already $C^{2,\alpha'}$.
The trace theorem then bounds its required
$W^{1-1/p,p}$ norm and explains the last term in
\eqref{g:W2p-preabsorption}.

The existing positive H\"older modulus makes the quadratic term
absorbable.  On a ball or boundary half-ball of radius $h_0$, subtract
a constant from $Z$.  First derivative interpolation, including the
scaled lower-order term \cite{Nirenberg1959}, gives
\begin{equation}\label{g:quadratic-interpolation}
 \|D_xZ\|_{2p}^2
 \leq C\operatorname{osc}Z\,\|D_x^2Z\|_p
       +Ch_0^{3/p-2}(\operatorname{osc}Z)^2.
\end{equation}
No homogeneous normal boundary condition is needed for this
interpolation.  In a flattened half-patch, choose
$c\in[\inf Z,\sup Z]$ and extend $u=Z-c$ by
$Eu(x',-t)=3u(x',t)-2u(x',2t)$ on a smaller collar and cut off on a
larger patch.  Its value and first normal derivative match across the
face, its sup norm is at most $5\operatorname{osc}Z$, and its
$W^{2,p}$ norm has the usual bounded extension estimate after scaling.
Whole-space interpolation yields
Equation~\eqref{g:quadratic-interpolation}, including its lower-order
scaled term.  The fixed smooth charts change only the constants.
The oscillation on each such patch is at most $Ch_0^\theta$.
Cover $U'$ by these smaller patches with bounded overlap, using a
slightly enlarged patch for each interpolation estimate.  Sum their
estimates in $\ell^p$.  This bounds the quadratic term by
$Ch_0^\theta\|D_x^2Z\|_p+C(h_0)$ on a slight enlargement.
Ordinary interpolation treats the $W^{1,p}$ term in the same fashion.
Take the supremum of the resulting $W^{2,p}$ norms over patches of one
fixed size throughout the truncation.  Each enlargement is covered by
a bounded number of other patches, independently of the truncation
radius.  This supremum is finite for each smooth solution on its finite
truncation.  Choose $h_0$ once so that the coefficient of the supremum
on the right is less than one half and absorb it.  We obtain a uniform
$W^{2,p}$ bound for $Z$ and hence a uniform $C^{1,\theta'}$ bound for
some $\theta'>0$.

Schauder estimates for the trace equation and then the expanded second
equation, with their respective boundary conditions, now increase
regularity successively.  At each step the source of the second
equation uses at most second derivatives of $f$, which have already
been controlled by the preceding trace estimate.  Iteration reaches
every required finite order, with constants depending on that order
and the fixed parameters.  The construction uses smooth
three-dimensional patches also at the rotational axis, so no singular
orbit-coordinate estimate is required.
\end{proof}
